\begin{lstlisting}[
    caption={Função utilizada para medir o tempo de execução},
    label={lst:calcular-tempo}
]
double calcularTempo(void (*algoritmo)(int*, int), int* vetor, int tamanho) {
    // salvar tempo do inicio
    auto inicio = chrono::steady_clock::now();

    // executar o algoritmo
    algoritmo(vetor, tamanho);

    // salvar o tempo do fim da execucao
    auto fim = chrono::steady_clock::now();

    // calcular duracao
    chrono::duration<double> tempo = fim - inicio;

    return tempo.count();
}
\end{lstlisting}

A função apresentada no Código~\ref{lst:calcular-tempo} é responsável por
medir o tempo de execução do algoritmo. Para isso, registra o instante
inicial, executa o algoritmo, registra o instante final e calcula a
diferença entre os dois instantes.

O valor retornado pela função representa o tempo de execução em segundos.